\section{Conclusion}
\label{sec:conclusion}

This thesis presented MalWhere, a pipeline that combines static and
dynamic malware analysis into confidence-scored, standards-compliant
threat intelligence: MITRE ATT\&CK technique mappings exported as
STIX~2.1 bundles and live MISP events. Two mechanisms extend it beyond
a conventional static-plus-dynamic design: a confidence reconciliation
model that recognizes agreement both across sources and across
technique specificity levels (\S\ref{sec:reconciliation}), and a
resubmission loop that correctly attributes a multi-stage sample's
distributed capability to the component that actually exhibits it
rather than collapsing it into the parent binary's own score
(\S\ref{sec:resubmission}). Positioned against the tools surveyed in
\S\ref{sec:related-work}, the pipeline is not a replacement for a
broad reputation aggregator such as VirusTotal, which no single-team
research pipeline can match for raw submission-corpus scale; it is the
step that comes after, turning that first-pass signal into a
reconciled, per-technique, evidence-linked mapping for the sample that
actually needs one.

Validated against three malware families with materially different
architectures (a single-stage ransomware binary, a feature-rich .NET
infostealer, and a four-component loader/rootkit/RAT chain) against
ground truth built from function-level manual reverse engineering, the
pipeline reached family-level F1 scores of 0.96, 0.89, and 0.84
respectively, with high precision throughout
(\S\ref{sec:evaluation}). Those numbers are only as trustworthy as the
process behind them, which is why this thesis places equal weight on
the audit trail that produced them: 15 discrepancies between pipeline
output and ground truth were individually traced to a specific,
checkable cause, not accepted or dismissed on the aggregate score
alone (\S\ref{sec:fp-audit}); a real aggregation bug that had silently
substituted the wrong analyzed binary was found and fixed
(\S\ref{sec:case-installer}); and static detection coverage was
extended from roughly 30 to approximately 85 MITRE technique
identifiers, every addition sourced from MITRE's own technique
descriptions and confirmed not to introduce new findings on the
validated set before being trusted (\S\ref{sec:coverage}).

The central argument of this work is that this audit trail, every rule
tied to verifiable evidence, every discrepancy diagnosed to a specific
cause, every extension checked for regression before being kept, is a
more defensible contribution than the headline F1 numbers by
themselves, and is the property most likely to transfer when the
pipeline meets a malware family it has not been tuned against.
Precisely because that transfer has not yet been measured, expanding
the validated family set (\S\ref{sec:future-work}) is identified as
the single highest-value next step, alongside the concretely-scoped
extension to Linux/ELF-based IoT botnet analysis once its two missing
prerequisites, real ELF-aware detection rules and a ground-truth Linux
sample, are in place.
